% Source: https://artofproblemsolving.com/community/c3h2949413_2019_emcc_guts_round_p13_p24__exeter_math_club_competition
% Stable identity: aops-topic-2949413-q23
Let the parallel lines be $y=0$ and $y=2r$, with the third centre $(0,r)$. The first centre has height 1, and the second height $2r-2$, since both lie in the strip. External tangency to the third gives their horizontal coordinates $\pm2\sqrt r$ and $\pm2\sqrt{2r}$. Their own tangency requires horizontal separation squared $3^2-(2r-3)^2=4r(3-r)$. If their horizontal signs are opposite, the squared separation is $4r(3+2\sqrt2)$, forcing negative $r$, impossible. With the same sign it is $4r(3-2\sqrt2)$, hence $r=\boxed{2\sqrt2}$. This configuration fits the strip and meets every required tangency.
\begin{asy}
usepackage("amsmath");
size(7cm);
real r=2*sqrt(2);pair O=(0,r),A=(2*sqrt(r),1),B=(2*sqrt(2*r),2*r-2);draw(circle(O,r),blue);draw(circle(A,1));draw(circle(B,2));draw((-3,0)--(8,0));draw((-3,2*r)--(8,2*r));label("$r$",O);label("$1$",A);label("$2$",B);label("$M$",(-3,0),W);label("$N$",(-3,2*r),W);
\end{asy}
